\documentclass[11pt]{article}

\usepackage[margin=1.05in]{geometry}
\usepackage{amsmath,amssymb,amsthm,mathtools}
\usepackage{thmtools}
\usepackage{enumitem}
\usepackage[hidelinks]{hyperref}
\usepackage{microtype}

\newtheorem{theorem}{Theorem}[section]
\newtheorem{lemma}[theorem]{Lemma}
\newtheorem{proposition}[theorem]{Proposition}
\newtheorem{corollary}[theorem]{Corollary}
\newtheorem{definition}[theorem]{Definition}
\newtheorem{remark}[theorem]{Remark}

\newcommand{\Z}{\mathbb Z}
\newcommand{\N}{\mathbb N}
\newcommand{\R}{\mathbb R}
\newcommand{\T}{\mathbb T}
\newcommand{\GI}{\mathbb Z[i]}
\newcommand{\Norm}[1]{\left\lVert#1\right\rVert^2}
\newcommand{\floor}[1]{\left\lfloor #1 \right\rfloor}
\newcommand{\ceil}[1]{\left\lceil #1 \right\rceil}
\newcommand{\conj}[1]{\overline{#1}}
\newcommand{\E}{\mathcal E}
\newcommand{\M}{\mathcal M}

\title{A Sub-Logarithmic Bound for Lattice Points on Small Circular Arcs}
\author{Craig Chen}
\date{}

\begin{document}
\maketitle

\begin{abstract}
    Let $I$ be an arc of the circle of radius $R$ centered at the origin of length at most $C\sqrt{R}$. We prove that 
$$
\#(I \cap \Z^2) \ll_C \frac{\log R}{\log\log R}
$$
for every $C>0$.
The proof combines a Ramana-type determinant identity over \(\Z[i]\) and leverages residue collisions modulo primes not dividing $R$ to establish divisibility of the determinant, which we can translate to a lower bound on the spacing between points on the arc. 
%Divisibility is not guaranteed a priori, but Mertens' first theorem guarantees an interval of primes with enough weighted mass to run a dichotomy: either enough of its primes miss $R$ and the collision divisibility fires, or some prime divides $R$ and split/inert descent in $\Z[i]$ passes the family to a smaller circle."
The result is formalized in Lean v4.30.
\end{abstract}

%\noindent\textbf{Keywords.}
%lattice points on circles; short arcs; Gaussian integers; determinant method;
%restriction estimates; Lean 4.

\tableofcontents

\section{Introduction}
\label{sec:intro}

Let $\E_R \coloneq \{ (x,y) \in \mathbb{Z}^2 \, | \, x^2 + y^2 = R^2 \}$ be the set of lattice points on a centered circle of radius $R$.
We will always assume $R^2 \in \mathbb{N}$.
The problem can be approached from two different ways.
Arithmetically, \(\E_R\) is controlled by the factorization of \(R^2\) in the Gaussian integers.
Geometrically, once we zoom into a very short arc, clustering is a statement about distances: one can examine the chord lengths between lattice points.
We will frequently move back and forth between these two perspectives.

Previous results are able to prove finite bounds for arcs of length $R^{\alpha}$ when $\alpha \in [0,1/2)$ but only yield a $\log R$ bound at the critical exponent $\alpha=1/2$.
The scale \(\sqrt R\) is indeed critical for the circle: an arc of this size is just long enough that it has nontrivial sagitta.
It is conjectured that arcs of length comparable to \(R^{1/2}\) also contain only finitely many lattice points, uniformly in the radius \(R\), \cite{CillerueloGranville2007,CillerueloGranville2009}.
For almost all \(R\) there are quite few lattice points on such arcs \cite{BourgainRudnick2011}; it is uniformity in \(R\) that remains hard (see Section~\ref{subsec:almost-all}).
In this paper we improve the existing logarithmic bound by proving a sublogarithmic estimate that every arc \(I\) of length at most \(C\sqrt R\) satisfies
\[
  \#(I\cap\E_R)\ \ll_C\ \frac{\log R}{\log\log R}.
\]
The proof is a determinant method strengthened by extra arithmetic divisibility.
A Vandermonde-type determinant over \(\Z[i]\) cannot vanish when it consists of distinct points on the circle, but if many consecutive points somehow lie in a short interval then the same determinant is forced to be small.
This is the obstruction to clustering we leverage to arrive at our result.
The logarithmic saving comes from strengthening the determinant with congruences modulo rational primes.
Primes \(p \nmid R^2\) give collision divisibility, while primes dividing \(R^2\) allow descent to a smaller circle.
Balancing these two alternatives on a single interval of primes, whose weighted mass Mertens' first theorem controls, produces the \(\log\log R\) gain.

\subsection{Main result}
\label{subsec:main-result}

We identify a point \((x,y)\in\Z^2\) with \(z=x+iy\in\GI\), and write \(\Norm{z}=z\conj z=x^2+y^2\).

\begin{restatable}[Geometric arc form]{theorem}{geometricthm}\label{thm:geometric}
For every \(C>0\) and \(K>8\) there exists \(R_0(C,K)\) such that the following holds.
For all \(R\ge R_0(C,K)\) with \(R^2\in\N\), let \(I\) be an arc of the circle \(x^2+y^2=R^2\) of length at most \(C\sqrt R\).
Then
\[
  \#(I\cap\E_R)
  \le
  K\,\frac{\log R}{\log\log R}.
\]
\end{restatable}

%This is a radius statement.  It does not require \(R\) itself to be an integer;
%the arithmetic input is only that the squared radius \(R^2\) is the integer
%norm \(N\).  Equivalently, one may set \(R=\sqrt N\) for an integer norm \(N\),
%in which case the length scale is \(C N^{1/4}\) and the logarithmic factor is
%\(\log\sqrt N/\log\log\sqrt N\).

\subsection{Proof sketch}
\label{subsec:proof-idea}

We now give a more explicit roadmap of the proof; precise definitions are provided in Section \ref{sec:mainresult}.
The key parameter in the argument is the block size $s$.
This is ultimately where the $\log / \log\log$ shape comes from. 


The first step in the proof is to notice that a determinant identity over \(\GI\) turns clustering into an integrality problem: from any \(2s+1\) distinct points of \(\E_R\) it produces a nonzero Gaussian integer that is forced to be small when all the pairwise chords are short.
Given points \(z_0,\dots,z_{2s}\in\E_R\), let \(\mathcal R_s(z)\) be the \((2s+1)\times(2s+1)\) matrix whose \(j\)-th column is
\[
  \bigl(\conj{z_j}^{\,s},\ \conj{z_j}^{\,s-1},\ \ldots,\ \conj{z_j},\ 1,\ z_j,\ \ldots,\ z_j^{\,s}\bigr)^T,
\]
and set \(J_s=\det\mathcal R_s(z)\), a Gaussian integer.
This is the Vandermonde matrix of the Laurent monomials \(z^{-s},\dots,z^s\) in disguise: on the circle \(\conj z=R^2/z\), so multiplying the \(j\)-th column by \(z_j^s\) clears every conjugate, and the powers of \(R^2\) created along the way factor out row by row, leaving the ordinary Vandermonde matrix in the \(z_j\).
Carrying this out (Proposition~\ref{prop:ramana}) gives the Ramana-type identity \cite{Ramana2007}
\[
  \Bigl(\prod_{j=0}^{2s} z_j^s\Bigr)\,J_s
  =
  R^{s(s+1)}\prod_{0\le i<j\le 2s}(z_j-z_i).
\]
Since the points are distinct, the Vandermonde product on the right is nonzero, and the prefactor is nonzero because \(R>0\); hence \(J_s\ne0\).
Now take norms throughout the identity.
Each point satisfies \(\Norm{z_j}=R^2\), so the prefactor contributes \(\prod_j\lVert z_j\rVert^{2s}=R^{2s(2s+1)}\) and the norm of the Vandermonde product is the product of all squared chord lengths.
Altogether
\[
  R^{2s(2s+1)}\,\Norm{J_s}
  =
  R^{2s(s+1)}\prod_{i<j}\Norm{z_j-z_i},
\]
and dividing through, since \(s(2s+1)-s(s+1)=s^2\),
\[
  \Norm{J_s} = \frac{\prod_{i<j} \Norm{z_j - z_i}}{R^{2s^2}}\,:
\]
the norm of the determinant is the product of all squared chord lengths of the block, measured in units of \(R^{2s^2}\).
Because \(J_s\) is a nonzero Gaussian integer, \(\Norm{J_s}\ge1\), which is to say the product of the squared chord lengths is at least \(R^{2s^2}\).
In other words, \(2s+1\) points on the circle cannot all be pairwise close: their chords must be long on average, with geometric mean at least \(R^{1/2-1/(4s+2)}\).
This is the obstruction to clustering.
Quantitatively, a block of \(2s+1\) points with all pairwise chords at most \(B\) has all squared chord lengths at most \(B^2\), hence
\[
  R^{2s^2}\le B^{2s(2s+1)},
\]
so whenever \(B^{2s(2s+1)}<R^{2s^2}\), no such block can exist: among any \(2s+1\) of the points, some pair is more than \(B\) apart.

The determinant is a fixed-size tool: it takes exactly \(2s+1\) points and says that those particular points cannot all be pairwise close.
Feeding all \(M\) points into one giant determinant doesn't help because enlarging the block grows the chord budget and the norm requirement simultaneously.
If the obstruction held at the full arc length \(B=L\), nothing more would be needed: any \(2s+1\) of the points are pairwise within \(L\), so \(M\le 2s\) outright.
But it doesn't.
On its own, the obstruction is effective only below the scale \(B\approx R^{1/2-1/(4s+2)}\), well short of \(L\approx C\sqrt R\).
Applied directly at that scale it still gives a bound, by cutting the arc into \(L/B\) stretches of length \(B\), but the covering factor \(L/B\) produces the old logarithmic bound.
To improve the result we need to instead improve the determinant lower bound \(\Norm{J_s}\ge1\) until the geometric obstruction holds at the full arc length.

The second main step is to notice that the determinant must abide by certain divisibility constraints.
Take some prime \(p\nmid R^2\).
Modulo such a prime, the equation \(x^2+y^2=R^2\) has only \(\ll p\) possible residue classes.
A block of \(2s+1\) points therefore has many pairs which collide modulo \(p\).
Each collision contributes a factor of \(p\) to the determinant.
If many primes in a controlled interval are missing from \(R^2\), the accumulated prime-power divisibility makes the determinant too divisible to be compatible with all chords being small.
This is what we call the "missing-prime" section of the proof.

We are not guaranteed to find many primes missing from $R^2$ as needed above, though, which leads us to the third step of the proof.
If the missing-prime branch fails, then many primes in the same interval must divide \(R^2\).
In particular, one finds a prime divisor \(p\) of \(R^2\) in a controlled range.
The arithmetic of \(\Z[i]\) then gives descent.
If \(p\equiv 3\pmod 4\), the rational prime is inert and divides every point on the circle when it divides the norm.
If \(p\equiv 1\pmod 4\), the prime splits as \(p=\pi\overline\pi\), and at least half of the points are divisible by one of \(\pi,\overline\pi\).
Dividing by the relevant factor maps the chosen points to a smaller circle and shrinks chord lengths.
The determinant obstruction can then be applied on the smaller circle.

The final proof balances these two alternatives on a single interval of primes: either sufficiently many primes are missing, giving extra determinant divisibility, or a useful prime divisor of \(R^2\) appears, allowing us to reduce the problem to a smaller circle.
Choosing the parameter \(s\) at scale \(\log R/\log\log R\) makes the two mechanisms meet, and this is where the sublogarithmic bound comes from.

\subsubsection*{Why \(\log R/\log\log R\)?}

Where does \(\log R/\log\log R\) come from, and why can't we simply choose something finite?

Taking logarithms in the block inequality \(R^{2s^2}\le B^{2s(2s+1)}\), a block of \(2s+1\) points with pairwise chords at most \(B\) is impossible if
\[
  s(2s+1)\log B^2<2s^2\log R,
  \qquad\text{i.e.}\qquad
  B<R^{\frac12-\frac{1}{4s+2}}.
\]
But what we want is \(B\asymp\sqrt R\).

For context, here is what a direct application gives without any improvement of the determinant lower bound.
We apply the block obstruction at the largest scale it can forbid, \(B_s=R^{1/2-1/(4s+2)}\), and cut the arc into
\[
  \frac{L}{B_s}\ \le\ C\,\frac{R^{1/2}}{R^{1/2-1/(4s+2)}}\ =\ C\,R^{\frac{1}{4s+2}}
\]
stretches of length \(B_s\), each holding at most \(2s\) points, so
\[
  M\ \lesssim\ s+s\,C\,R^{\frac{1}{4s+2}},
\]
which is minimized at \(s=\tfrac14\log R\).
This is the same logarithmic result that we can get from the Cilleruelo--C\'ordoba argument.

We can do better than just $\log R$ because the missing-prime argument allows us to improve the determinant lower bound from \(\Norm{J_s}\ge1\) to \(\Norm{J_s}\ge\prod_p p^{2c_p}\), where \(c_p\) counts congruence collisions in the block.
Extra divisibility raises the size of $B_s$, but to increase \(B_s=R^{1/2-1/(4s+2)}\) all the way up to the critical scale \(C\sqrt R\), the gain must be at least $s\log R$.
By counting prime collisions in the block and using Mertens' on the total gain we get to our $\log\log R$ bound.
%Each missing prime \(p\) yields about \(s^2/2p\) collisions in the window --- pigeonhole against the \(\le2p\) residue classes of the mod-\(p\) circle --- and each collision is worth \(2\log p\) of divisibility, so the prime \(p\) contributes a gain of about \(s^2\log p/p\).
%Summing, the total gain is \(s^2\) times the weighted mass \(\sum\log p/p\) of the missing primes, and Mertens' theorem evaluates such sums: \(\sum_{p\le x}\log p/p\approx\log x\).
%The catch is that only small primes can collide at all --- the window must overfill the \(\le2p\) residue classes, so only \(p\lesssim s\) contributes (the proof works with an interval \((m_0,m_1]\) of primes with \(4m_1\le s\)).
%The available mass is therefore at most \(\sum_{p\le s}\log p/p\approx\log s\), and the largest achievable gain \(s^2\log s\) meets the deficit \(s\log R\) only if
%\[
%  s\log s\ \gtrsim\ \log R,
%  \qquad\text{i.e.}\qquad
%  s\ \gtrsim\ \frac{\log R}{\log\log R}.
%\]
Since every branch of the proof returns \(M\ll s\), we take \(s\) as small as this balance allows, which is \(s\asymp\log R/\log\log R\). 

%The same accounting answers the second question: no fixed window size can work.
%For constant \(s\), the total obstruction a window can accumulate --- the norm gap plus all available collision divisibility --- is at most about \(s^2\log s\), a quantity independent of \(R\), while the deficit \(s\log R\) grows without bound.
%The strengthened window inequality therefore fails for every fixed \(s\) once \(R\) is large: within this framework, arcs of length \(C\sqrt R\) cannot be given a uniformly finite bound.
%Finally, the missing primes are only conditionally available, since every prime in \((m_0,m_1]\) might instead divide \(R^2\).
%The proof handles this with a dichotomy: a Chebyshev bound guarantees that each interval \((m_08^i,m_08^{i+1}]\) carries weighted prime mass at least \(\log2/2\), so across \(k\asymp\log\log R\) such intervals either the missing mass reaches the scale \(\log\log R\) demanded above, or some prime of controlled size divides \(R^2\) and descent applies, again yielding \(M\ll s\).

\medskip
\noindent\textbf{Overview.}
In Section~\ref{sec:related} we review classical results in this problem and compare the proof present to that of Cilleruelo and C\'ordoba.
Section~\ref{sec:mainresult} states the main result and all the ingredients of the proof.
Section~\ref{sec:formalization} describes the Lean formalization.
Full proofs are in the appendices.
A glossary of the notation used throughout the paper is provided in Appendix~\ref{sec:glossary}.

\section{Related Work}
\label{sec:related}

\subsection{Convex-curve bounds and determinant methods}
\label{subsec:determinant-methods}

The classical benchmark is Jarn\'{\i}k's theorem: a strictly convex arc of length \(L\) passes through \(\ll L^{2/3}\) lattice points \cite{Jarnik1926}.
Applied directly to a circular arc of length \(L=C\sqrt R\), this gives \(\ll_C R^{1/3}\) which takes advantage of the curvature of the arc but does not leverage any structure of the circle.
Bombieri and Pila's determinant method gives far-reaching bounds for integral points on arcs and ovals \cite{BombieriPila1989}; our proof is also determinant-driven, but it uses the special arithmetic of a fixed circle and the Gaussian integers.
The identity at its core is the Gaussian-integer case of a Vandermonde identity of Ramana \cite{Ramana2007}, introduced there with applications like ours in mind: integer points on small arcs of conics and divisors of an integer in prescribed congruence classes.

The determinant method also has a non-archimedean branch which resembles our divisibility step.
Heath-Brown's \(p\)-adic determinant method \cite{HeathBrown2002} replaces the archimedean input of Bombieri--Pila with a congruence: points that agree modulo an auxiliary prime \(p\) make the determinant divisible by a large power of \(p\), just as points that are close on a short arc make it small.
Salberger's global determinant method \cite{Salberger2023} accumulates the divisibility over many primes at once, from the way the points distribute over the residue classes of the curve modulo each prime, and Walsh \cite{Walsh2015} sharpened the bookkeeping for curves.

Written in this language, the missing-prime mechanism of Section~\ref{sec:mainresult} is the Vandermonde case of the global method, run on a fixed circle.
The arc supplies the archimedean smallness, residue collisions at the missing primes supply the divisibility, and the nonzero integer \(\Norm{J_s}\) has to satisfy both constraints at once, yielding the inequality \eqref{eq:missing-prime-condition}; the collision counts \(c_p\) play the role of the class statistics of the global method.
%Our auxiliary primes are tiny --- below the block size, where pigeonhole forces collisions across the roughly \(p\) residue classes --- whereas Heath-Brown takes \(p\) comparable to a power of the height and works one class at a time.
Primes dividing \(R^2\), which the general method has no use for, get a separate mechanism here: descent in \(\GI\), which is special to norm forms.
And the circle offers a structural simplification: the determinant never vanishes (Proposition~\ref{prop:ramana}), so the divisibility cannot be absorbed by a vanishing determinant and instead forces up the norm lower bound.
Could this machinery strengthen the theorem?
Not the shape, as far as we can tell.
The deeper \(p\)-adic filtration --- collisions modulo \(p^2,p^3,\ldots\) --- is a convergent constant gain, the pigeonhole cap of Section~\ref{sec:future} binds the global method exactly as it binds ours, and running the method at a single large prime reproduces the below-critical fixed-\(k\) bounds rather than improving the critical scale.
We view the present proof as the small-prime, collision-based corner of the global determinant method, with descent as the extra arithmetic input.

\subsection{Cilleruelo--C\'ordoba and Cilleruelo--Granville}
\label{subsec:cc-log}

For circles centered at the origin, much sharper local results are known.
Cilleruelo and C\'ordoba proved fixed-\(k\) bounds below the square-root scale \cite{CillerueloCordoba1992}.
Namely: for every \(k\ge 1\), every arc of length
\[
  \sqrt2\,R^{1/2-\frac{1}{4\floor{k/2}+2}}
\]
contains at most \(k\) lattice points.
Cilleruelo and Granville later studied the fine structure of close lattice points on circles and gave families of four points at short scales \cite{CillerueloGranville2009}; related conjectures on boundedness of short arcs appear in \cite{CillerueloGranville2007}.

The Cilleruelo--C\'ordoba estimates imply a logarithmic bound at the \(\sqrt R\) scale.
Assume \(R>e^2\), and set
\[
  k_R=\ceil{\frac{\log R}{2}}.
\]
Then
\[
  4\floor{k_R/2}+2\ge \log R
  \implies
  \frac{1}{4\floor{k_R/2}+2}\le \frac1{\log R}.
\]
%Indeed, if \(k_R=2q\), then \(2q\ge \frac12\log R\), so
%\(4q+2\ge \log R+2\).  If \(k_R=2q+1\), then
%\(2q+1\ge \frac12\log R\), so \(4q+2\ge \log R\). 
The length allowed by the Cilleruelo--C\'ordoba theorem is at least
\[
  \sqrt2\,R^{1/2-\frac{1}{4\floor{k_R/2}+2}}
  \ge
  \sqrt2\,R^{1/2-\frac1{\log R}}
  =
  \frac{\sqrt2}{e}\sqrt R.
\]
Thus arcs of length at most \((\sqrt2/e)\sqrt R\) contain at most
\[
  k_R\le \frac12\log R+1
\]
lattice points.

For a general fixed \(C>0\), divide an arc of length at most \(C\sqrt R\) into
\[
  J_C=\ceil{\frac{Ce}{\sqrt2}}
\]
subarcs, each of length at most \((\sqrt2/e)\sqrt R\).
Summing the preceding bound over the subarcs gives
\[
  \#(I\cap\E_R)
  \le
  J_C\left(\frac12\log R+1\right)
  \ll_C
  \log R.
\]

Both our approach and Cilleruelo--C\'ordoba's approach use integrality-gap arguments in \(\GI\) by isolating a quantity that is a nonzero Gaussian integer (hence has norm at least \(1\)) and that is forced to be small when many lattice points crowd into a short arc.
In their argument \cite{CillerueloCordoba1992} they study the differences in angles of the lattice poitnts.
Writing the split part of \(R^2\) as \(\prod_j p_j^{\alpha_j}\), every point of \(\E_R\) has the form \(R\,e^{2\pi i(\sum_j\gamma_j\Phi_j+t/4)}\), where \(\Phi_j\) is the angle of the Gaussian prime above \(p_j\) and the integer exponents satisfy \(|\gamma_j|\le\alpha_j\) and \(\gamma_j\equiv\alpha_j\pmod2\).
The half-angle difference of two points is then, up to a unit, the angle of a lattice point on a circle of radius at most \(\sqrt2\prod_j p_j^{|\gamma_j-\gamma_j'|/4}\); being off the real axis, that point has imaginary part at least \(1\), so the angular gap between the two points is at least the reciprocal of this smaller radius.
Multiplying the pair bound over all \(\binom{m+1}2\) pairs and maximizing the exponent sums \(\sum_{s,s'}|\gamma_{j,s}-\gamma_{j,s'}|\) by an extremal combinatorial argument yields the fixed-\(k\) theorem.

Notice that for a single block the two mechanisms are quantitatively identical (Corollary~\ref{cor:subcritical}).
The Ramana determinant should therefore be viewed as a repackaging of the same integrality gap: the Vandermonde factorization collects all pairwise differences of a block into a single nonzero Gaussian integer, and the norm bound \(\Norm{J_s}\ge1\) is applied once \textit{per block} instead of once \textit{per pair}.
Repackaging this way lets us use primes outside the factorization of \(R^2\).
The determinant, being a polynomial in the coordinates of the points, also feels congruences from primes that do not divide $R^2$.
The main difference between the approaches is that theirs relies on the prime divisors of $R^2$ whereas ours relies on the non-divisors.
It's this second, unconditional supply of primes that improves the bound from \(\log R\) to \(\log R/\log\log R\).

\subsection{The logarithmic barrier}
\label{subsec:log-barrier}

Could the logarithm of the previous subsection be an artifact of the argument?
So far it seems like no: the exponent \(1/2-1/(4\floor{k/2}+2)\) has not been improved for any \(k\) since 1992, and for arcs of length \(C\sqrt R\) nothing between \(\ll_C\log R\) and the conjectured \(\ll_C1\) is known.

\emph{Arcs of length \(R^{\alpha}\), \(\alpha<1/2\).}
The exponent \(1/2-1/(4\floor{k/2}+2)\) is sharp for \(k\le3\), including the constants.
An arc of length \(\sqrt2\) contains at most one lattice point, and the pairs \((n,n+1)\), \((n+1,n)\) span arcs of length \(\sqrt2+o(1)\) \cite{CillerueloGranville2007}.
Cilleruelo proved that an arc of length \((16R)^{1/3}\) contains at most two lattice points \cite{Cilleruelo1991}, and the triples
\[
  (4n^3-1,\ 2n^2+2n),\qquad (4n^3,\ 2n^2+1),\qquad (4n^3+1,\ 2n^2-2n)
\]
realize that length asymptotically.
Cilleruelo and Granville determined the sharp constant for three points at the same exponent \(R^{1/3}\), and constructed infinite families of circles with four lattice points on arcs of that length; one of the families is built from Fibonacci numbers \cite{CillerueloGranville2009}.
For every \(k\ge4\) the 1992 theorem is still the best known.
Already for \(k=4\) the question is open in both directions: there is no construction of infinitely many circles with four lattice points on an arc of length \(o(R^{2/5})\), and no result rules them out \cite{CillerueloGranville2007}.

\emph{Arcs of length \(C\sqrt R\).}
Cilleruelo and Granville also stated the boundedness conjecture in coordinates: for every \(\alpha<1/2\), the number of representations \(R^2=a^2+b^2\) with \(|b|\) in a prescribed interval of length \(R^{\alpha}\) should be bounded in terms of \(\alpha\) alone \cite[Conjecture~13]{CillerueloGranville2007}.
They observe that intervals starting at \(b=0\) are elementary for \(\alpha\le1/2\): if \(|b|\le\sqrt R\), then \(a=\sqrt{R^2-b^2}\) is confined to an interval of length at most \(1\).
In arc language: an arc of length \(C\sqrt R\) that touches one of the four points where the circle meets the axes contains finitely lattice points for trivial reasons.
The difficulty is uniformity over the position of the arc, and the interval version is open for every \(\alpha>1/2\).

\emph{Arcs of length \(R^{\alpha}\), \(\alpha>1/2\).}
Here the best known uniform bound is the divisor bound \(\#\E_R=R^{o(1)}\) for the entire circle.
Cilleruelo and Granville conjecture uniform boundedness all the way up to length \(R^{1-\varepsilon}\) \cite[Conjecture~14]{CillerueloGranville2007}.

%\emph{Determinant methods.}
%The determinant methods of Bombieri--Pila, Heath-Brown, Salberger, and Walsh \cite{BombieriPila1989,HeathBrown2002,Salberger2023,Walsh2015} bound rational points of bounded height on curves, uniformly in the curve, and progress there is measured in the exponent of the height.
%On the fixed circle \(x^2+y^2=R^2\) they prove nothing sharper than the divisor bound, and in particular cannot distinguish \(\log R\) from a constant.

We have found no bound between \(\ll_C\log R\) and the conjectured finite bound in the literature; Theorem~\ref{thm:geometric} appears to be the first sublogarithmic bound for arcs of length \(C\sqrt R\).

\subsection{Squares in arithmetic progressions and a congruence analogue}
\label{subsec:squares-ap}

Why is the scale \(\sqrt R\) so challenging?
One answer is that a uniform bound for arcs of length \(C\sqrt R\) is connected to old open problems about squares.
Let \(\sigma(k)\) denote the maximal number of squares in an arithmetic progression of length \(k\).
Rudin conjectured \(\sigma(k)\ll\sqrt k\) \cite{Rudin1960}.
Szemer\'edi proved \(\sigma(k)=o(k)\) \cite{Szemeredi1974}, Bombieri, Granville and Pintz proved \(\sigma(k)\ll k^{2/3+o(1)}\) \cite{BombieriGranvillePintz1992}, and Bombieri and Zannier proved \(\sigma(k)\ll k^{3/5+o(1)}\) \cite{BombieriZannier2002}, which is still the record; the last two proofs pass through counting rational points on families of auxiliary curves.
Cilleruelo and Granville connect these problems to the clustering of lattice points on circles through a list of intermediate conjectures \cite{CillerueloGranville2007}.

Another parallel concerns the congruence analogue of a short arc: how many solutions can \(x^2\equiv a\pmod b\) have in an interval of length \(\sqrt b\)?
Konyagin and Steger proved that a degree-\(d\) congruence has \(\ll 1+\log L/\log(1+b^{1/d}/L)\) roots in any interval of length \(L\), which at \(L=b^{1/d}\) is \(\ll\log b\) \cite{KonyaginSteger1994}.
For \(d=2\), Cilleruelo and Granville rederive the logarithmic bound by a Vandermonde argument: the product \(\prod_{i<j}(x_j-x_i)\) over solutions in a short interval is a nonzero integer divisible by a large power of \(b\), and comparing its size against its divisibility bounds the number of solutions \cite{CillerueloGranville2007}.
They conjecture that the number of solutions in an interval of length \(\sqrt b\) is in fact bounded, and show that this implies Rudin's conjecture.

In both problems, every known upper bound at the square-root scale comes from such a comparison and produces a logarithm.
Theorem~\ref{thm:geometric} shows that on the circle the logarithm is not the limit of these arguments.
The theorem does not approach uniform boundedness, and it has no direct consequence for Rudin's conjecture.
%Does the missing-prime mechanism have an analogue mod \(b\), where \(\prod_{i<j}(x_j-x_i)\) picks up extra divisibility from small primes not dividing \(b\)?
%We find this an interesting question.

\subsection{Uniform versus almost-all radii}
\label{subsec:almost-all}

What makes this problem particularly challenging is that the "average" circle and corresponding arcs contain very few lattice points.
It's important to distinguish the uniform problem from its almost-all counterpart because the results are completely different.
Bourgain and Rudnick observed that outside a sparse exceptional set of radii, every pair of distinct lattice points on the circle of radius $R$ is separated by at least $R^{1-\varepsilon}$ \cite{BourgainRudnick2011}.
In other words, for almost all \(R\) an arc of length \(R^{1-\varepsilon}\) contains at most \emph{one} lattice point.
On the statistical side, Kurlberg and Lester proved that along a density-one subsequence of admissible radii the nearest-neighbor spacings of the angles of \(\E_R\) are Poissonian, and also that the density-one restriction is unavoidable: there are subsequences, even with \(\#\E_R\to\infty\), whose angles cluster so strongly that the spacing measure degenerates to a point mass at zero \cite{KurlbergLester2025}.
Clustering is an exceptional phenomenon, and much of what we know about the worst cases comes from computation.

There's also related work that restricts to a special subset of radii. Chan obtains logarithmic improvements beyond the \(\sqrt R\) scale for \(R^2\) of the form a square plus a much smaller square \cite{Chan2015}.
Temur studies the concentration of lattice points on short arcs of general conics through sieving and diophantine approximation, obtaining uniform bounds for weighted counts rather than for the counts themselves \cite{Temur2020}.
However, neither approach yields a uniform sublogarithmic bound for arcs of length \(C\sqrt R\) for general \(R\).

%\subsection{Restriction estimates}

%The same local-clustering question enters restriction estimates for toral
%eigenfunctions.  If
%\[
%  \varphi(x)=\sum_{\mu\in\E_R}a_\mu e^{2\pi i\langle \mu,x\rangle}
%\]
%is an eigenfunction on \(\T^2\), then restriction of \(\varphi\) to a curved
%hypersurface is controlled by oscillatory sums over pairs of frequencies.
%Near-parallel frequencies, which are precisely frequencies lying in short arcs
%of \(\E_R\), form the difficult local contribution.  Bourgain and Rudnick
%proved uniform restriction bounds for curves with nonzero curvature
%\cite{BourgainRudnick2009}.  The present theorem gives a formalized
%sublogarithmic local multiplicity estimate at the critical scale, and therefore
%supplies a quantitative input for any restriction argument whose loss is
%controlled by such arc-clustering numbers.

\section{Main Results}
\label{sec:mainresult}

Full proofs are deferred to the appendix.
We first present the key components.
Throughout this section, \(C>0\) is fixed, and \(I\) is an arc of the circle \(x^2+y^2=R^2\) of length \(L\le C\sqrt R\).
We write \(z_0,\dots,z_{M-1}\) for the lattice points of \(I\); chords are bounded by arclength, so any two of them satisfy \(|z_i-z_j|\le L\).

\subsection{The Ramana determinant}
\label{subsec:ramana}

Let \(s\ge 1\), and let \(z_0,\dots,z_{2s}\in\GI\).
Define the \((2s+1)\times(2s+1)\) matrix
\[
  \mathcal R_s(z)_{r,j} =
  \begin{cases}
    \conj{z_j}^{\,s-r}, & 0\le r<s,\\
    1, & r=s,\\
    z_j^{\,r-s}, & s<r\le 2s.
  \end{cases}
\]
Let
\[
  J_s(z_0,\ldots,z_{2s})=\det \mathcal R_s(z).
\]
The following identity is the case \(A=\GI\), \(\alpha_j=\conj{z_j}\), \(\beta_j=z_j\), \(\gamma=R^2\) of Theorem~1 of Ramana \cite{Ramana2007}.

\begin{restatable}[Ramana determinant identity]{proposition}{ramanaidentity}\label{prop:ramana}
If \(z_j\conj z_j=R^2\) for every \(j\), then
\[
  \left(\prod_{j=0}^{2s}z_j^s\right)J_s(z_0,\ldots,z_{2s})
  =
  R^{s(s+1)}\prod_{0\le i<j\le 2s}(z_j-z_i).
\]
If the \(z_j\) are distinct, then \(J_s\ne 0\).
\end{restatable}

\begin{restatable}[Subcritical block obstruction]{corollary}{subcriticalobstruction}\label{cor:subcritical}
Let \(z_0,\dots,z_{2s}\in\GI\) be distinct and satisfy \(z_j\conj z_j=R^2>0\).
If the chord bound \(|z_i-z_j|\le B\) holds for every \(i<j\), then
\[
  R^{2s^2}\le \floor{B^2}^{s(2s+1)}.
\]
Consequently, if \(\floor{B^2}^{s(2s+1)}<R^{2s^2}\), some pair in the block has \(|z_i-z_j|>B\).
\end{restatable}

On an arc shorter than the forbidden scale, the whole family is a single block: if \(B\ge L\) and \(\floor{B^2}^{s(2s+1)}<R^{2s^2}\), then any \(2s+1\) of the points would have all pairwise chords at most \(L\le B\), so \(M\le 2s\).
We call this the \emph{saturated} regime.
Run at a scale \(B<L\), the obstruction still bounds \(M\), at the price of cutting the arc into \(L/B\) stretches of length \(B\); that covering computation is the baseline bound of Section~\ref{subsec:proof-idea}.
The proof is arranged so that the covering is never needed: every branch runs its obstruction at a scale equal to its arc length.

\subsection{Missing primes}
\label{subsec:missing-primes}

For interval endpoints \(m_0<m_1\), let
\[
  \mathcal P_R(m_0,m_1)=\{p\text{ prime}:m_0<p\le m_1,\ p\nmid R^2\}
\]
be the set of missing primes in the interval \((m_0,m_1]\), and define the weighted missing-prime mass
\[
  \M_R(m_0,m_1)
  =
  \sum_{p\in\mathcal P_R(m_0,m_1)}\frac{\log p}{p}.
\]

Fix a block \(z_0,\dots,z_{2s}\) of distinct points of \(\E_R\).
The missing-prime mechanism strengthens the lower bound \(\Norm{J_s}\ge1\) used in Corollary~\ref{cor:subcritical} in three steps: residue collisions modulo each missing prime force divisibility of the chord products, the divisibility survives cancellation into the determinant, and the accumulated coprime prime powers raise the norm lower bound.

\begin{restatable}[Residue collisions]{lemma}{residuecollisions}\label{lem:collision}
Let \(p\) be prime.
The congruence \(x^2+y^2\equiv R^2\pmod p\) has at most \(2p\) solutions in \((\Z/p\Z)^2\).
Consequently, if
\[
  c_p=\#\{(i,j):0\le i<j\le 2s,\ z_i\equiv z_j \pmod p\},
\]
where the congruence means both coordinates agree modulo \(p\), then
\[
  c_p\ \ge\ \frac{(2s+1)^2}{4p}-\frac{2s+1}{2}.
\]
If moreover \(4p\le s\), then \(c_p\ge s^2/(2p)\).
\end{restatable}

\begin{restatable}[Determinant divisibility]{proposition}{determinantdivisibility}\label{prop:divisibility}
Let \(p\nmid R^2\) be prime.
Then \(p^{c_p}\mid J_s(z_0,\ldots,z_{2s})\) in \(\GI\), and hence \(p^{2c_p}\mid\Norm{J_s}\) in \(\Z\).
If \(P\) is any finite set of such primes, then
\[
  \prod_{p\in P}p^{2c_p}\ \le\ \Norm{J_s}.
\]
\end{restatable}

\begin{restatable}[Enhanced block obstruction]{corollary}{enhancedobstruction}\label{cor:enhanced}
Let \(z_0,\dots,z_{2s}\in\GI\) be distinct with \(z_j\conj z_j=R^2>0\), and let \(P\) be a finite set of primes not dividing \(R^2\).
If the chord bound \(|z_i-z_j|\le B\) holds for every \(i<j\), then
\[
  R^{2s^2}\prod_{p\in P}p^{2c_p}\ \le\ \floor{B^2}^{s(2s+1)}.
\]
\end{restatable}

Taking logarithms, Corollary~\ref{cor:enhanced} says that a block of \(2s+1\) points in which every chord length is at most \(B\) is impossible if
\[
  \underbrace{s(2s+1)\log \floor{B^2}}_{\text{chord budget}}
  \;<\;
  \underbrace{2s^2\log R}_{\text{norm gain}}
  \;+\;
  \underbrace{\sum_{p\in P}2c_p\log p}_{\text{collision gain}}.
\]
The left side is the total logarithmic size of all \(s(2s+1)\) chord norms, each a rational integer of size at most \(B^2\); the first term on the right is just the Ramana obstruction of Corollary~\ref{cor:subcritical}; the second is the missing-prime enhancement.
Now replace each quantity by the explicit bound supplied by the preceding results:
\begin{itemize}[leftmargin=2em]
\item for \(P=\mathcal P_R(m_0,m_1)\) with \(4m_1\le s\), Lemma~\ref{lem:collision} gives \(c_p\ge s^2/(2p)\ge E_p\coloneq\floor{s^2/(2p)}\);
\item rounding each exponent to an integer loses at most one collision, i.e.\ at most \(2\log p\) of gain: \(2\floor{s^2/(2p)}\ge s^2/p-2\), whence
\[
  \sum_{p\in P}2E_p\log p
  \;\ge\;
  s^2\sum_{p\in P}\frac{\log p}{p}-2\sum_{p\in P}\log p
  \;=\;
  s^2\M_R(m_0,m_1)-2\sum_{p\in P}\log p .
\]
\end{itemize}
A missing prime \(p\) produces about \(s^2/(2p)\) collisions in the block, each worth \(2\log p\) of divisibility, so the gain contributed by \(p\) per unit of \(s^2\) is exactly its Mertens mass \(\log p/p\).
Substituting in the chord budget yields the finite weighted-log condition
\[
  s(2s+1)\log \floor{B^2}
  <
  s^2\M_R(m_0,m_1)-2\sum_{\substack{m_0<p\le m_1\\p\nmid R^2}}\log p
  +2s^2\log R ,
  \tag{MP}\label{eq:missing-prime-condition}
\]
which is thus nothing more than the logarithmic form of Corollary~\ref{cor:enhanced} with every block-dependent quantity replaced by its explicit bound.
%The exact Lean statement uses natural-number floors and coercions; this display is its mathematical content.

\begin{remark}[Reading \eqref{eq:missing-prime-condition} at the critical
scale]\label{rem:mp-critical} Substitute \(B=C\sqrt R\), so that \(\log\floor{B^2}\le\log R+2\log C\), and divide by \(s^2\).
Then \eqref{eq:missing-prime-condition} is implied by
\[
  \M_R(m_0,m_1)
  \;>\;
  \frac{\log R}{s}+4\log C+o(1),
\]
where the \(o(1)\) collects the floor toll \(2\sum_{p\in P}\log p/s^2\ll(\log m_1)/s\) and a \((2\log C)/s\) term; up to these lower-order terms, \eqref{eq:missing-prime-condition} also requires this.
So the condition says precisely that the weighted missing-prime mass must beat the per-\(s^2\) deficit \((\log R)/s\) identified in Section~\ref{subsec:proof-idea}.
With \(s\asymp\log R/\log\log R\) the demand is of size \(\log\log R\), which is what the Mertens interval of the next subsection supplies.
\end{remark}

\begin{restatable}[Missing-prime branch]{proposition}{missingprimebranch}\label{prop:missing}
Let \(s\ge1\), let \(m_0<m_1\) with \(4m_1\le s\), and let \(B>0\) with \(\floor{B^2}\ge1\).
If \eqref{eq:missing-prime-condition} holds, then no \(2s+1\) of the points have all pairwise chords at most \(B\).
In particular, if \(B\ge L\), then \(M\le 2s\).
\end{restatable}

\subsection{Prime extraction and descent}
\label{subsec:descent}

The missing-prime condition \eqref{eq:missing-prime-condition} is only useful if missing primes exist, and it's not clear if that's the case a priori in a given range.
The proof needs a supply of primes that is guaranteed to be present at finite, explicit sizes.
Conveniently, Mertens' first theorem provides exactly this \cite{Mertens1874}: for every \(x\ge1\),
\[
  -2-E_1\ \le\ \sum_{p\le x}\frac{\log p}{p}-\log x\ \le\ \log 4+4,
  \qquad
  E_1\coloneq\sum_{p}\frac{\log p}{p(p-1)}\le\frac{5\log2+3}{4}.
\]
Let \(2\le m_0\le m_1\).
Evaluating at both endpoints and subtracting gives a floor for the weighted mass of the interval:
\[
  \M(m_0,m_1)\coloneq\sum_{m_0<p\le m_1}\frac{\log p}{p}
  \ \ge\ \log\frac{m_1}{m_0}-c_M,
  \qquad
  c_M\coloneq(\log4+4)+(2+E_1)\le 10 .
\]
The deficit \(c_M\) is a fixed constant, so taking \(m_1/m_0\) equal to a power of \(\log R\) makes the floor a multiple of \(\log\log R\).

The proof branches into two paths based on the following condition.
The missing-prime mechanism needs primes non-divisors with enough weighted mass, while descent needs a prime divisor of \(R^2\).
Neither resource is unconditionally available.
What we can do is choose a threshold that always guarantees at least one of the branches will have enough resources.
We call this the \emph{few-primes condition}, which holds when the missing primes account for less than the guaranteed floor:
\[
  \M_R(m_0,m_1)<\log\frac{m_1}{m_0}-c_M.
  \tag{FP}\label{eq:few-primes}
\]
This says that measured in weighted mass, few of the interval's primes are missing from the factorization of \(R^2\).
We call this branching the \emph{few-primes dichotomy}.
If \eqref{eq:few-primes} is not true, then the missing-prime mass is at least \(\log(m_1/m_0)-c_M\); with \(\log(m_1/m_0)\asymp\log\log R\) this is exactly the gain of size \(\log\log R\) that Remark~\ref{rem:mp-critical} demands for the missing-prime branch.
If \eqref{eq:few-primes} holds, then since the full mass of \((m_0,m_1]\) is at least the floor, some prime of the interval is not missing: there is a prime divisor
\[
  m_0<p\le m_1,\qquad p\mid R^2,
\]
for the descent approach to use.

\begin{restatable}[Descent branch]{proposition}{descentbranch}\label{prop:descent}
Let \(p\mid R^2\) be a prime with \(m_0<p\le m_1\).
\begin{enumerate}[label=(\alph*)]
\item If \(p\equiv 3\pmod 4\), every \(z\in\E_R\) is divisible by \(p\) in \(\GI\), and division by \(p\) maps the family into \(\E_{R/p}\).
Chord lengths are divided by \(p\).
\item If \(p\equiv 1\pmod 4\), write \(p=\pi\conj\pi\) in \(\GI\).
At least half the points are divisible by one of \(\pi,\conj\pi\).
Dividing that subfamily by the corresponding Gaussian prime maps it into \(\E_{R/\sqrt p}\).
Chord lengths are divided by \(\sqrt p\).
\end{enumerate}
If the corresponding descended subcritical inequalities hold, the original family satisfies \(M\ll s\).
\end{restatable}

\subsection{Scales and branch bounds}
\label{subsec:scales}

The few-primes dichotomy leaves two situations: either the missing mass on \((m_0,m_1]\) reaches the Mertens floor, or some prime \(m_0<p\le m_1\) divides \(R^2\).
Each one feeds the block obstruction, and the only choice left is the scale \(B\) at which to run it.
We take each scale equal to the largest arc length its branch can see: the missing branch sees the original arc, and descent shortens the arc by a factor of \(\sqrt p\) or \(p\):
\[
  B=C\sqrt R,
  \qquad
  B_{\mathrm{split}}(p)=\frac{C\sqrt R}{\sqrt p},
  \qquad
  B_{\mathrm{inert}}(p)=\frac{C\sqrt R}{p}.
\]
At these scales the obstruction is saturated: it forbids \(2s+1\) points within the full arc length, so it applies to the whole family at once.

\begin{restatable}[Branch bounds]{proposition}{branchbounds}\label{prop:branches}
Let \(s\ge1\) and \(2\le m_0<m_1\) with \(4m_1\le s\).
\begin{enumerate}[label=(\alph*)]
\item If \eqref{eq:missing-prime-condition} holds at \(B=C\sqrt R\), then \(M\le 2s\).
\item If some prime \(p\mid R^2\) with \(m_0<p\le m_1\) and \(p\equiv3\pmod4\) satisfies
\[
  s(2s+1)\log\floor{B_{\mathrm{inert}}(p)^2}<s^2\log\frac{R^2}{p^2},
\]
then \(M\le 2s\).
\item If some prime \(p\mid R^2\) with \(m_0<p\le m_1\) and \(p\equiv1\pmod4\) satisfies
\[
  s(2s+1)\log\floor{B_{\mathrm{split}}(p)^2}<s^2\log\frac{R^2}{p},
\]
then \(M\le 4s\).
\end{enumerate}
\end{restatable}

The displayed inequalities are the subcritical block obstruction of Corollary~\ref{cor:subcritical} on the descended circles \(\E_{R/p}\) and \(\E_{R/\sqrt p}\).
The split branch is the only one that pays a factor of \(2\), for the divisibility class of Proposition~\ref{prop:descent}.

It remains to choose the parameters, and this choice is where the shape \(\log R/\log\log R\) appears.
Given \(K>8\), put \(Q=K/4>2\), fix \(q\) and \(\theta_1\) with \(2<q<Q\) and \(\tfrac12+\tfrac1q<\theta_1<1\) (possible since \(Q>2\)), and take
\[
  s=\ceil{q\,\frac{\log R}{\log\log R}},\qquad
  m_0=\floor{\sqrt{\log R}},\qquad
  m_1=\floor{(\log R)^{\theta_1}}.
\]
Once \(R\) is large, these choices satisfy \(4m_1\le s\) and the inequality of Proposition~\ref{prop:branches} in whichever branch the few-primes dichotomy hands us: the missing branch demands \((1+o(1))\,(\log\log R)/q\) of weighted mass against the Mertens floor \((\theta_1-\tfrac12-o(1))\log\log R\), and the descent inequalities reduce to \((1+o(1))\,(\log\log R)/q<\log p\) with \(\log p>(\tfrac12-o(1))\log\log R\).
The appendix carries out the verification.
In either case
\[
  M\le 4s\le 4\,Q\,\frac{\log R}{\log\log R}=K\,\frac{\log R}{\log\log R},
\]
which proves Theorem~\ref{thm:geometric}.
The endpoint of this schedule is optimal for the argument.
With \(m_0=(\log R)^{\theta}\) the missing branch needs \(Q>1/(\theta_1-\theta)\) and the descent endpoints need \(Q>1/\theta\), and \(\max(1/\theta,\,1/(\theta_1-\theta))\ge2\) for every \(\theta<\theta_1<1\), with equality approached exactly at \(\theta=1/2\), \(\theta_1\to1\).

\section{Formalization}
\label{sec:formalization}

Every step of the proof, from the determinant identity through Mertens' first theorem with explicit constants to the final parameter choice, is machine-checked in Lean~4 over Mathlib \cite{Mathlib2020,deMouraUllrich2021}; the Mertens input is adapted from the PrimeNumberTheoremAnd project \cite{PNTplus}.
The formalized statement is the following parametrized form of Theorem~\ref{thm:geometric}.
It avoids a topological formalization of arcs: in place of an arc it takes ordered real parameters whose differences dominate chord lengths, with arclength along the arc as the motivating instance.
The integer norm appears as its own variable \(N\), with \(R^2=N\).

\begin{restatable}[Asymptotic parametrized theorem]{theorem}{paramthm}\label{thm:param}
Let \(C>0\) and \(K>8\).
There exists \(R_0(C,K)\) such that for every \(R\ge R_0(C,K)\), the following holds.
Let \(M,N\in\N\), \(a,L\in\R\), and let \(z:\N\to\GI\), \(t:\N\to\R\).
Assume
\[
  R^2=N,\qquad 0<M,\qquad 0<N,\qquad L\le C\sqrt R,
\]
\[
  z_i\conj{z_i}=N\quad\text{for all }i<M,
\]
that \(z_i\ne z_j\) for distinct \(i,j<M\), that \(t_i\le t_j\) whenever \(i\le j<M\), that
\[
  |z_i-z_j|^2\le (t_j-t_i)^2\qquad (i,j<M),
\]
and that
\[
  a\le t_i\le a+L\qquad (i<M).
\]
Then
\[
  M\le K\,\frac{\log R}{\log\log R}.
\]
\end{restatable}

This is the statement of \texttt{mertens\_asymptotic\_parametrized\_theorem} in \texttt{MertensMain.lean}, proved for every \(K>8\), exactly as stated; the proof is the one in the appendix.
What is machine-checked end to end is Theorem~\ref{thm:param}.
Theorem~\ref{thm:geometric} follows by taking \(t_i\) to be arclength; this reduction is the elementary chord--arclength estimate, and we write it out at the start of Appendix~\ref{app:asymptotic}.
The arclength wrapper\\
\texttt{eventually\_jarnik\_arclength\_circular\_arc\_sublog\_of\_constant\_gt\_eight}\\
in \texttt{ArcToParameter.lean} states Theorem~\ref{thm:geometric} in this form: it takes the chord--arclength inequality as a hypothesis rather than deriving it from a formalized theory of arcs, so the reduction lives on paper.

The development is organized as follows.

\begin{itemize}[leftmargin=2em]
\item \texttt{RamanaDeterminant.lean}: the determinant identity and nonvanishing.
\item \texttt{SubcriticalBound.lean}: the subcritical block obstruction.
\item \texttt{ModPrimeCollision.lean} and \texttt{DeterminantDivisibility.lean}: collisions modulo missing primes and determinant divisibility.
\item \texttt{MissingPrimeIntervalBranch.lean}: the weighted missing-prime branch.
\item \texttt{MertensFirst.lean} and \texttt{MertensIntervalMass.lean}: Mertens' first theorem with explicit constants and the interval mass floor.
\item \texttt{SaturatedWindows.lean}: the saturated block bounds behind Proposition~\ref{prop:branches}.
\item Gaussian descent:
\begin{itemize}[nosep,leftmargin=1.5em]
\item \texttt{PrimeDescent.lean}
\item \texttt{DescentGeometry.lean}
\item \texttt{DescentSubcritical.lean}
\end{itemize}
\item \texttt{MertensDichotomy.lean}: the few-primes dichotomy and the branch bounds of Proposition~\ref{prop:branches}.
\item \texttt{MertensSchedule.lean} and \texttt{MertensMain.lean}: the asymptotic parameter choice and the final theorem.
\end{itemize}

The project builds with \texttt{lake build} without any uses of \texttt{sorry}, \texttt{admit}, or project-level \texttt{axiom} declarations.
The development is available at \url{https://github.com/craigxchen/Jarnik}; it pins Lean \texttt{v4.30.0} and a fixed Mathlib commit (\texttt{c5ea003}) through its manifest.

Appendix~\ref{sec:lean-correspondence} maps each numbered statement of the paper to the Lean declaration that formalizes it.
%\section{Numerical Search Code}
%
%The repository also contains Rust code for experimental searches for unusually
%dense endpoint-scale arcs.  This code is not used as an assumption in the Lean
%proof; it is a way to find and inspect candidate families, especially possible
%five- and six-point clusters.
%
%The optimized command is
%\[
%\texttt{cargo run --release -- arc-families --n-max N --min-size 5}.
%\]
%It differs from the full diagnostic command \(\texttt{search}\) in three ways.
%First, it builds a smallest-prime-factor sieve up to \(N\), so each integer in
%the scan can be factored by repeated table lookup rather than by trial division.
%Second, it skips any norm whose factorization cannot produce enough lattice
%representations to meet the requested family size.  Third, during the scan it
%computes only the maximum endpoint-scale cluster size; it materializes and
%serializes a compact JSONL record only after the threshold is met.
%
%The emitted record contains the norm \(N\), the radius \(R=\sqrt N\), the
%angular endpoint window, the actual angular width, the lattice points in the
%cluster, and the split-prime data for the norm.  Thus a typical five- or
%six-point search can be run as a lightweight discovery pass, and selected norms
%can then be passed to \(\texttt{analyze-n}\) for the full cut, relation-height,
%transition-pattern, and half-angle diagnostics.

%\section{Conclusion and Restriction Consequences}
%
%To recap, the theorem gives an improvement from the Cilleruelo--C\'ordoba logarithmic baseline \(\ll_C\log R\) to
%\[
%  \#(I\cap\E_R)\ \ll_C\ \frac{\log R}{\log\log R}
%\]
%for arcs \(I\) of length \(C\sqrt R\).
%
%The proof stays entirely inside the integrality-gap framework of Section~\ref{subsec:cc-log}.
%What is new is the bookkeeping --- a determinant that can absorb congruence information at primes not dividing \(R^2\), and a dichotomy that converts the absence of such primes into descent to a smaller circle.
%
%The result also supplies a quantitative input for restriction estimates on \(\T^2\).
%The frequency set of an eigenfunction with eigenvalue \(4\pi^2R^2\) is exactly \(\E_R\), and the difficult local contribution to restriction kernels comes from pairs of nearly parallel frequencies, that is, frequencies lying in short arcs of \(\E_R\).
%Any restriction argument whose local loss is controlled by the maximal number of lattice points in arcs of length \(C\sqrt R\) may therefore replace a logarithmic loss by \(\log R/\log\log R\), in a machine-checked form.
%We emphasize that this is an arithmetic input to such arguments, not a standalone restriction theorem: converting the cluster bound into an operator-norm estimate requires the analytic framework, for instance that of Bourgain and Rudnick \cite{BourgainRudnick2009}, in which the cluster number appears.

%For restriction problems on \(\T^2\), the frequency set of an eigenfunction is
%\(\E_R\).  Restriction to a curved hypersurface leads to kernels of the form
%\[
%  K(s,t)=\sum_{\mu,\nu\in\E_R}
%  a_\mu\overline{a_\nu}\,
%  e^{2\pi i\langle \mu-\nu,\gamma(s)-\gamma(t)\rangle}.
%\]
%Pairs \((\mu,\nu)\) with nearly parallel directions are the part least affected
%by curvature and stationary phase.  The number of such frequencies is controlled
%by short-arc multiplicity.  Therefore any restriction argument whose local loss
%is bounded by the maximal number of lattice points in arcs of length
%\(C\sqrt R\) may replace a logarithmic loss by
%\[
%  \frac{\log R}{\log\log R}.
%\]
%In this sense Theorem~\ref{thm:geometric} supplies a formal arithmetic input for
%restriction estimates: it quantifies, in a machine-checked way, that critical
%frequency clusters are sublogarithmic.  It should not be read as a standalone
%new restriction theorem; translating the local cluster estimate into a sharp
%operator norm requires the analytic restriction framework in which the cluster
%number appears.

\section{Future Work}
\label{sec:future}

Can the method of this paper be pushed all the way to the conjectured uniform bound \cite{CillerueloGranville2007}, under which arcs of length \(C\sqrt R\) contain only finitely many lattice points?
It seems unlikely.
The constant in the bound can probably be optimized a bit further, but the $\log R / \log\log R$ shape is unlikely to change because of two things.

First, only small primes can contribute.
A block of \(2s+1\) points has collisions modulo \(p\) only when it overfills the roughly \(p\) residue classes of the mod-\(p\) circle, which requires \(p\lesssim s\).
Even if we count the collisions exactly, the total divisibility a block can collect is at most about \(s^2\sum_{p\le s}\log p/p\approx s^2\log s\).
Second, the chord bound \(\floor{B^2}\) looks like the most wasteful step of the proof because it bounds everything with the maximum possible distance but sharpening it does not help.
Even if we could use the exact distances, we would only get an improvement (after taking logarithms) of size $O(s^2)$.
Both gains are \(\ll s^2\log s\), while reaching the critical scale requires \(s\log R\) (Section~\ref{subsec:proof-idea}).
A finite bound will have to come from a different kind of argument.

What could a different kind of argument look like?
Both limits above apply only to proofs that take a local view of clusters of lattice points. 
One candidate is the structure theory of close lattice points of Cilleruelo and Granville \cite{CillerueloGranville2009}. 
Another is a hybrid of determinant divisibility with the multiplicative parameterization of \cite{CillerueloCordoba1992}.
Their angular method has sharp per-pair bounds but aggregates them lossily because the configurations that saturate its combinatorial step never lie on one short arc, so congruence information might rule them out.
Finally, the descent branch consumes only a single prime divisor of \(R^2\) at a time, while the conjectured obstruction should strengthen with every split prime.
Using several divisors at once seems within reach of the present framework.

\appendix

\section{Proofs}
\label{sec:proofs}

\subsection{The determinant and subcritical bounds}
\label{app:determinant}

\ramanaidentity*

\begin{proof}
The \(j\)-th column of \(\mathcal R_s\) is
\[
  (\conj z_j^s,\conj z_j^{s-1},\ldots,\conj z_j,1,z_j,\ldots,z_j^s)^T.
\]
Multiply this column by \(z_j^s\).
Because \(z_j\conj z_j=R^2\), the entries become
\[
  (R^{2s},R^{2(s-1)}z_j,\ldots,R^2z_j^{s-1},z_j^s,z_j^{s+1},\ldots,z_j^{2s})^T.
\]
Factor \(R^{2(s-r)}\) from the row corresponding to \(\conj z^{s-r}\), for \(r=0,\dots,s-1\).
The total factored power is
\[
  R^{2(1+2+\cdots+s)}=R^{s(s+1)}.
\]
The remaining matrix is the ordinary Vandermonde matrix with columns
\[
  (1,z_j,z_j^2,\ldots,z_j^{2s})^T.
\]
Its determinant is
\[
  \prod_{0\le i<j\le 2s}(z_j-z_i).
\]
This proves the identity.
If the \(z_j\) are distinct, the Vandermonde product is nonzero in the domain \(\GI\), and the product \(\prod_j z_j^s\) is nonzero because \(R>0\).
Hence \(J_s\ne0\).
\end{proof}

\subcriticalobstruction*

\begin{proof}
Take norms in Proposition~\ref{prop:ramana}.
Since \(J_s\) is a nonzero Gaussian integer,
\[
  \Norm{J_s}\ge 1.
\]
The norm of \(\prod_j z_j^s\) is
\[
  \prod_j \left(\Norm{z_j}\right)^s=R^{2s(2s+1)}.
\]
The norm of the right side is
\[
    \Norm{ R^{s(s+1)} }
  \prod_{i<j}\Norm{z_j-z_i}
  =
  R^{2s(s+1)}\prod_{i<j}\Norm{z_j-z_i},
\]
because the rational integer \(R^{s(s+1)}\) has Gaussian norm \(R^{2s(s+1)}\).
Thus
\[
  R^{2s(2s+1)}
  \le
  R^{2s(s+1)}\prod_{i<j}\Norm{z_j-z_i},
\]
and hence
\[
  R^{2s^2}\le \prod_{i<j}\Norm{z_j-z_i}.
\]
There are
\[
  \binom{2s+1}{2}=s(2s+1)
\]
pairs.
Each norm is a rational integer, so if every chord length is at most \(B\), then every norm is at most \(\floor{B^2}\) and
\[
  R^{2s^2}\le \floor{B^2}^{s(2s+1)}.
\]
The final assertion is the contrapositive.
\end{proof}

\subsection{The missing-prime and descent mechanisms}
\label{app:missing}

\residuecollisions*

\begin{proof}
For each of the \(p\) values of \(x\in\Z/p\Z\), the congruence \(y^2\equiv R^2-x^2\pmod p\) is a quadratic equation over the field \(\Z/p\Z\) and has at most two roots \(y\).
Hence the mod-\(p\) circle has at most \(2p\) points.

Write \(n=2s+1\) and let \(\rho(z)=(\Re z\bmod p,\ \Im z\bmod p)\) denote coordinatewise reduction.
Since \(\Re(z_i)^2+\Im(z_i)^2=R^2\) in \(\Z\), the map \(\rho\) sends the block into the mod-\(p\) circle.
For each residue class \(b\) let \(f_b\) be the number of indices \(i\) with \(\rho(z_i)=b\), so that \(\sum_b f_b=n\) with at most \(2p\) nonzero terms.
By the Cauchy--Schwarz inequality,
\[
  n^2=\Bigl(\sum_b f_b\Bigr)^{\!2}\le 2p\sum_b f_b^2 .
\]
The number of ordered pairs \((i,j)\) with \(\rho(z_i)=\rho(z_j)\), the diagonal included, is \(\sum_b f_b^2\); removing the \(n\) diagonal pairs and halving,
\[
  c_p=\frac{\sum_b f_b^2-n}{2}\ \ge\ \frac{n^2}{4p}-\frac n2 .
\]
For the second claim, the inequality \(n^2/(4p)-n/2\ge s^2/(2p)\) is equivalent to \(n^2-2pn\ge 2s^2\).
From \(4p\le s\),
\[
  n^2-2pn\ \ge\ (2s+1)^2-\tfrac s2(2s+1)\ =\ 3s^2+\tfrac72 s+1\ \ge\ 2s^2 .
  \qedhere
\]
\end{proof}

\determinantdivisibility*

\begin{proof}
If \(z_i\equiv z_j\pmod p\), both coordinates of \(z_j-z_i\) are divisible by \(p\), so \(p\mid z_j-z_i\) in \(\GI\).
The \(c_p\) colliding pairs index distinct factors of the Vandermonde product, whence
\[
  p^{c_p}\ \Big|\ \prod_{0\le i<j\le 2s}(z_j-z_i).
\]
By Proposition~\ref{prop:ramana},
\[
  \Bigl(\prod_j z_j^s\Bigr) J_s
  =
  R^{s(s+1)}\prod_{i<j}(z_j-z_i),
\]
so \(p^{c_p}\) divides \(u\,J_s\), where \(u=\prod_j z_j^s\).

We now cancel \(u\).
Its norm is \(\Norm{u}=R^{2s(2s+1)}\), which is coprime to \(p\) because \(p\nmid R^2\) and \(p\) is prime.
In general, if \(u,d\in\GI\) with \(p\nmid\Norm{u}\), and \(p^e\mid u\,d\) in \(\GI\), then \(p^e\mid d\).
Indeed, multiplying by \(\conj u\) gives \(p^e\mid \conj u\,u\,d=\Norm{u}\,d\); a rational integer divides a Gaussian integer exactly when it divides both coordinates, so \(p^e\) divides \(\Norm{u}\Re(d)\) and \(\Norm{u}\Im(d)\) in \(\Z\); since \(\gcd(p^e,\Norm{u})=1\), it divides \(\Re(d)\) and \(\Im(d)\), that is, \(p^e\mid d\) in \(\GI\).
Applying this with \(d=J_s\) and \(e=c_p\) gives \(p^{c_p}\mid J_s\).

Taking norms, \(\Norm{p^{c_p}}=p^{2c_p}\) divides \(\Norm{J_s}\).
For distinct primes the powers \(p^{2c_p}\) are pairwise coprime, so their product over \(P\) also divides \(\Norm{J_s}\).
Since the \(z_j\) are distinct, \(J_s\ne0\) by Proposition~\ref{prop:ramana}, so \(\Norm{J_s}\ge\prod_{p\in P}p^{2c_p}\).
\end{proof}

\enhancedobstruction*

\begin{proof}
As in the proof of Corollary~\ref{cor:subcritical}, the chord norms are rational integers at most \(\floor{B^2}\), so taking norms in the determinant identity gives
\[
  R^{2s(2s+1)}\,\Norm{J_s}
  =
  R^{2s(s+1)}\prod_{i<j}\Norm{z_j-z_i}
  \le
  R^{2s(s+1)}\floor{B^2}^{s(2s+1)},
\]
hence \(R^{2s^2}\Norm{J_s}\le \floor{B^2}^{s(2s+1)}\).
Inserting the lower bound of Proposition~\ref{prop:divisibility} for \(\Norm{J_s}\) completes the proof.
\end{proof}

\missingprimebranch*

\begin{proof}
Set \(P=\mathcal P_R(m_0,m_1)\) and \(E_p=\floor{s^2/(2p)}\).

Suppose some \(2s+1\) of the points have all pairwise chords at most \(B\), so the chord bound of Corollary~\ref{cor:enhanced} holds for every pair of the block.

Every \(p\in P\) satisfies \(4p\le 4m_1\le s\), so Lemma~\ref{lem:collision} gives \(c_p\ge s^2/(2p)\ge E_p\) for the collision counts of the block.
Corollary~\ref{cor:enhanced} applied to the block therefore forces
\[
  R^{2s^2}\prod_{p\in P}p^{2E_p}
  \ \le\
  R^{2s^2}\prod_{p\in P}p^{2c_p}
  \ \le\
  \floor{B^2}^{s(2s+1)} .
\]
All quantities are at least \(1\), so we may take logarithms; using \(2E_p=2\floor{s^2/(2p)}\ge s^2/p-2\),
\[
  s(2s+1)\log \floor{B^2}
  \ \ge\ 2s^2\log R+\sum_{p\in P}2E_p\log p
  \ \ge\ 2s^2\log R+s^2\M_R(m_0,m_1)-2\sum_{p\in P}\log p .
\]
This contradicts \eqref{eq:missing-prime-condition}.
Hence no \(2s+1\) of the points have all pairwise chords at most \(B\), which is the first claim.
If \(B\ge L\), any \(2s+1\) of the points would have pairwise chords at most \(L\le B\); hence \(M\le 2s\).
\end{proof}

\descentbranch*

\begin{proof}
If \(p\equiv3\pmod4\), then \(-1\) is not a square modulo \(p\).
From
\[
  x^2+y^2\equiv0\pmod p
\]
we obtain \(x\equiv y\equiv0\pmod p\).
Hence \(p\mid z\) in \(\GI\) for every point \(z\), and \(z/p\in\GI\) satisfies
\[
  \Norm{z/p}=R^2/p^2.
\]
Moreover
\[
  |z_i/p-z_j/p|=|z_i-z_j|/p.
\]
Thus the parameter length in the descended problem is divided by \(p\).

If \(p\equiv1\pmod4\), then \(p=\pi\conj\pi\) in \(\GI\).
Since
\[
  p\mid z\conj z,
\]
at least one of \(\pi,\conj\pi\) divides \(z\).
Partition the original family according to which Gaussian prime divides \(z\); one class has at least half the points.
Divide that class by the corresponding Gaussian prime.
The norm drops from \(R^2\) to \(R^2/p\), and chord lengths are divided by
\[
  |\pi|=\sqrt p.
\]

In either case, once the descended determinant inequality is subcritical, Corollary~\ref{cor:subcritical} applies on the descended circle, saturated at the descended arc length.
In the split case the passage to the larger of two classes costs a genuine factor of \(2\).
Proposition~\ref{prop:branches} records the outcome.
\end{proof}

\subsection{The branch bounds}
\label{app:branches}

\branchbounds*

\begin{proof}
(a) Since \(B=C\sqrt R\ge L\), Proposition~\ref{prop:missing} gives \(M\le 2s\).

(b) By Proposition~\ref{prop:descent}(a), division by \(p\) maps the family into \(\E_{R/p}\), and chords shrink by \(p\), so any \(2s+1\) of the descended points have all pairwise chords at most \(L/p\le B_{\mathrm{inert}}(p)\).
The hypothesis is the logarithmic form of the subcritical block obstruction of Corollary~\ref{cor:subcritical} on \(\E_{R/p}\) at the scale \(B_{\mathrm{inert}}(p)\), which forbids exactly such a block.
Hence \(M\le 2s\).

(c) By Proposition~\ref{prop:descent}(b), at least half the points descend to \(\E_{R/\sqrt p}\), with chords shrunk by \(\sqrt p\), so any \(2s+1\) of them have all pairwise chords at most \(L/\sqrt p\le B_{\mathrm{split}}(p)\).
As in (b), the hypothesis forbids such a block, so the subfamily has at most \(2s\) points, and \(M\le 2\cdot 2s=4s\).
\end{proof}

\subsection{The asymptotic parameter estimates}
\label{app:asymptotic}

Theorem~\ref{thm:param} is the parametrized statement of Section~\ref{sec:formalization}; it implies Theorem~\ref{thm:geometric} by arclength parametrization, and we first record that reduction.

\geometricthm*

\begin{proof}[Proof of Theorem~\ref{thm:geometric} from Theorem~\ref{thm:param}]
List the lattice points \(z_0,\dots,z_{M-1}\) of \(I\) in order along the arc, and let \(t_i\in[0,L]\) be the arclength position of \(z_i\) measured from an endpoint of \(I\), so that \(t_i\le t_j\) whenever \(i\le j\).
Two points of the circle at angular distance \(\delta\) have chord length \(2R\sin(\delta/2)\le R\delta\), and \(R\delta\) is exactly the arclength between them, so \(|z_i-z_j|\le|t_j-t_i|\) for all \(i,j\).
The family \((z_i,t_i)\) therefore satisfies the hypotheses of Theorem~\ref{thm:param} with \(N=R^2\) and \(a=0\), and the bound on \(M\) follows.
\end{proof}

\paramthm*

\begin{proof}
Given \(K>8\), put \(Q=K/4>2\), fix \(q\) and \(\theta_1\) with \(2<q<Q\) and \(\tfrac12+\tfrac1q<\theta_1<1\), and adopt the schedule of Section~\ref{subsec:scales}:
\[
  s=\ceil{q\,\frac{\log R}{\log\log R}},\qquad
  m_0=\floor{\sqrt{\log R}},\qquad
  m_1=\floor{(\log R)^{\theta_1}} .
\]
Eventually \(2\le m_0\le m_1\).
We verify the budget \(s\le Q\log R/\log\log R\) and the inequalities of Proposition~\ref{prop:branches}.

\emph{Budget.}  \(s\le q\,\frac{\log R}{\log\log R}+1\le Q\,\frac{\log R}{\log\log R}\) eventually, because \((Q-q)\frac{\log R}{\log\log R}\to\infty\).

\emph{Interval.}  \(m_1\le (\log R)^{\theta_1}\) with \(\theta_1<1\), while \(s\ge q\log R/\log\log R\).
Since \(\log\log R=o\bigl((\log R)^{1-\theta_1}\bigr)\), it follows that \(4m_1\le s\) for all large \(R\).
In particular \(m_1=o(s)\).

\emph{Endpoints.}  Dividing the split inequality of Proposition~\ref{prop:branches}(c)
\[
  s(2s+1)\log\left(\frac{C^2R}{p}\right)<s^2\,(2\log R-\log p)
\]
by \(s^2\), the leading terms \(2\log R-\log p\) cancel and it becomes
\[
  4\log C+\frac1s\Bigl(\log R-\log p+2\log C\Bigr)\ <\ \log p .
\]
The left side is \((\log\log R)/q+O_C(1)\), because \((\log R)/s\le(\log\log R)/q\).
The right side exceeds \(\log m_0=\log\floor{\sqrt{\log R}}=(1/2-o(1))\log\log R\), and \(1/q<1/2\), so the inequality holds eventually, uniformly over the primes \(p>m_0\).
The inert inequality of Proposition~\ref{prop:branches}(b) reduces in the same way to \((\log\log R)/q+O_C(1)<2\log p\), which is weaker.

\emph{Threshold.}  By Remark~\ref{rem:mp-critical}, condition \eqref{eq:missing-prime-condition} at the scale \(B=C\sqrt R\) is implied by
\[
  \M_R(m_0,m_1)\ >\ \frac{\log R}{s}+4\log C+o(1),
\]
where the \(o(1)\) absorbs the floor toll \(2\sum_{p\in P}\log p/s^2\ll m_1\log m_1/s^2\), which vanishes because \(m_1=o(s)\).
The right side is \((1+o(1))\,(\log\log R)/q\).
In the missing branch the Mertens floor supplies
\[
  \M_R(m_0,m_1)\ \ge\ \log\frac{m_1}{m_0}-c_M
  \ \ge\ \Bigl(\theta_1-\frac12-o(1)\Bigr)\log\log R,
\]
using \(\log m_1\ge\theta_1\log\log R-\log2\) and \(\log m_0\le\tfrac12\log\log R\); the fixed deficit \(c_M\) vanishes into the \(o(1)\).
Since \(1/q<\theta_1-1/2\), supply exceeds demand for all large \(R\), so \eqref{eq:missing-prime-condition} holds at \(B=C\sqrt R\) in the missing branch.

In either branch of the few-primes dichotomy, Proposition~\ref{prop:branches} therefore applies, and
\[
  M\le 4s\le 4\,Q\,\frac{\log R}{\log\log R}=K\,\frac{\log R}{\log\log R}.
  \qedhere
\]
\end{proof}

The formalized proof is this argument verbatim: given \(K>8\) it takes \(q=(2+Q)/2\) and \(\theta_1\) the midpoint of \((\tfrac12+\tfrac1q,\,1)\), and each \(o(1)\) is an explicit eventual bound against the resulting positive margins.

\section{Correspondence with the Lean Formalization}
\label{sec:lean-correspondence}

Every declaration below lives in the \texttt{GaussianChain} namespace, sub-namespaced by its file; we list the declaration name and the file that contains it.
The hypotheses \texttt{OnCircleUpTo M N z} and \texttt{InjectiveUpTo M z} appearing in the formal statements describe the \(M\) points being counted at a fixed radius: the theorems are asymptotic in \(R\), and past the threshold they apply uniformly to every such \(M\)-point family at that radius.
This formulation avoids the false requirement that an infinite sequence of distinct points lie on one fixed circle.

\begin{itemize}[leftmargin=2em]
\item Proposition~\ref{prop:ramana} (Ramana determinant identity):\\
\texttt{prod\_pow\_mul\_det\_ramana\_eq\_normPow\_mul\_vandermondeProduct}\\
and \texttt{det\_ramana\_ne\_zero\_of\_injective} in \texttt{RamanaDeterminant.lean}.
\item Corollary~\ref{cor:subcritical} (subcritical block obstruction):\\
\texttt{not\_injective\_of\_pairwise\_sqDist\_le\_of\_pow\_lt} in \texttt{SubcriticalBound.lean}.
\item The saturation step of Section~\ref{sec:mainresult} (any \(2s+1\) points of the arc are one block):\\
\texttt{card\_le\_of\_window\_span\_saturated} in \texttt{SaturatedWindows.lean}.
\item The Mertens interval floor of Section~\ref{sec:mainresult}:\\
\texttt{weightedPrimeInterval\_ge\_log\_sub\_deficit} in \texttt{MertensIntervalMass.lean},\\
built on \texttt{Mertens.sum\_log\_prime\_div\_eq\_log} in \texttt{MertensFirst.lean}.
\item Lemma~\ref{lem:collision} (residue collisions):\\
\texttt{zmod\_circle\_card\_le\_two\_mul} in \texttt{ModPrimeCollision.lean} and\\
\texttt{pairCollisionCount\_lower\_bound\_on\_circle} in \texttt{DeterminantDivisibility.lean}.
\item Proposition~\ref{prop:divisibility} (determinant divisibility):\\
\texttt{prod\_pow\_two\_mul\_pairCollisionCount\_dvd\_det\_ramana\_norm\_natAbs}\\
in \texttt{DeterminantDivisibility.lean}.
\item Corollary~\ref{cor:enhanced} (enhanced block obstruction):\\
\texttt{not\_injective\_of\_pair\_norm\_pow\_lt\_missing\_prime\_product}\\
in \texttt{DeterminantDivisibility.lean}.
\item Proposition~\ref{prop:missing} (missing-prime branch):\\
\texttt{missing\_prime\_branch\_saturated} in \texttt{MertensDichotomy.lean}.
\item Proposition~\ref{prop:descent} (descent branch):\\
\texttt{card\_le\_of\_prime\_divisor\_descent\_branch\_log\_twoScale\_saturated}\\
in \texttt{MertensDichotomy.lean}.
\item Proposition~\ref{prop:branches} (branch bounds, with the few-primes dichotomy):\\
\texttt{card\_le\_four\_mul\_s\_of\_mertens\_interval\_log} and its radius-scale form\\
\texttt{card\_le\_four\_mul\_log\_div\_loglog\_of\_mertens\_radius\_sq\_endpoint\_bounds}\\
in \texttt{MertensDichotomy.lean}.
\item Theorem~\ref{thm:param} (asymptotic parametrized theorem):\\
\texttt{mertens\_asymptotic\_parametrized\_theorem} in \texttt{MertensMain.lean},\\
proved for every \(K>8\), exactly as stated.
\item Theorem~\ref{thm:geometric} (geometric arc form):\\
\texttt{eventually\_jarnik\_arclength\_circular\_arc\_sublog\_of\_constant\_gt\_eight}\\
in \texttt{ArcToParameter.lean}, an arclength-parametrized wrapper for every \(K>8\).
\end{itemize}

The declaration names of the main statements match the names used in this paper; the supporting declarations keep Mathlib-style descriptive names.

\section{Glossary of Notation}
\label{sec:glossary}

For reference, we collect the notation used throughout the paper, grouped by the role it plays in the proof.

\subsection*{The problem data}

\begin{center}
\begin{tabular}{@{}lp{0.76\textwidth}@{}}
\(N\) & Section~\ref{sec:formalization} only: the integer norm, \(N=R^2\), kept as an explicit variable of the Lean statement \\
\(R\) & radius of the circle, with \(R^2\in\N\); all asymptotics are in \(R\to\infty\) \\
\(\E_R\) & the set of lattice points on the circle, \(\{(x,y)\in\Z^2 : x^2+y^2=R^2\}\) \\
\(I\) & the arc under consideration \\
\(C\) & arc-length constant: arcs have length at most \(C\sqrt R\) \\
\(K\) & the constant of the main theorem, \(M\le K\log R/\log\log R\); every \(K>8\) is admissible \\
\(R_0(C,K)\) & the threshold radius past which the theorem holds \\
\(M\) & the number of lattice points in the family being counted \\
\(z_i\) & the lattice points, viewed as Gaussian integers \(x+iy\) \\
\(t_i\) & Section~\ref{sec:formalization} only: the ordered real parameters of the Lean statement, with \(|z_i-z_j|\le|t_j-t_i|\) \\
\(a,L\) & \(L\) is the arc length, \(L\le C\sqrt R\); in Section~\ref{sec:formalization}, all \(t_i\in[a,a+L]\)
\end{tabular}
\end{center}

\subsection*{Gaussian-integer notation}

\begin{center}
\begin{tabular}{@{}lp{0.76\textwidth}@{}}
\(\GI\) & the Gaussian integers \\
\(\Norm{z}\) & the norm \(z\conj z=|z|^2\), simultaneously arithmetic size and squared Euclidean length \\
\(\conj z\) & complex conjugate; on the circle, \(\conj z=R^2/z\) \\
\(\pi,\conj\pi\) & conjugate Gaussian primes above a split rational prime \(p=\pi\conj\pi\)
\end{tabular}
\end{center}

\subsection*{Block and determinant machinery}

\begin{center}
\begin{tabular}{@{}lp{0.76\textwidth}@{}}
\(s\) & block-size parameter: the obstructions apply to blocks of \(2s+1\) points; chosen \(\asymp\log R/\log\log R\) in the final schedule \\
\(\mathcal R_s(z)\) & the \((2s+1)\times(2s+1)\) Ramana matrix \\
\(J_s\) & its determinant, the nonzero Gaussian integer carrying the obstruction \\
\(B\) & chord bound: no \(2s+1\) points may have all pairwise chords at most \(B\); chord norms in such a block are at most \(\floor{B^2}\); the branch values are \(B_{\mathrm{split}}(p)=C\sqrt R/\sqrt p\) and \(B_{\mathrm{inert}}(p)=C\sqrt R/p\) \\
\(c_p\) & the number of colliding pairs modulo \(p\) in a block \\
\(E_p\) & the explicit lower bound \(\floor{s^2/(2p)}\le c_p\) for the collision count, used in \eqref{eq:missing-prime-condition} \\
\(P\) & a finite set of primes not dividing \(R^2\), in practice \(\mathcal P_R(m_0,m_1)\)
\end{tabular}
\end{center}

\subsection*{Primes and the interval}

\begin{center}
\begin{tabular}{@{}lp{0.76\textwidth}@{}}
\(p\) & a rational prime; missing means \(p\nmid R^2\) \\
\(m_0,m_1\) & endpoints of the prime interval \((m_0,m_1]\); in the final schedule \(m_0=\floor{\sqrt{\log R}}\) and \(m_1=\floor{(\log R)^{\theta_1}}\) \\
\(\mathcal P_R(m_0,m_1)\) & the missing primes in \((m_0,m_1]\) \\
\(\M_R(m_0,m_1)\) & the weighted missing-prime mass \(\sum_{p\in\mathcal P_R}\log p/p\)
\end{tabular}
\end{center}

\subsection*{The asymptotic parameter schedule}

\begin{center}
\begin{tabular}{@{}lp{0.76\textwidth}@{}}
\(Q\) & budget constant, \(s\le Q\log R/\log\log R\), set to \(Q=K/4\) (Section~\ref{subsec:scales} and Appendix~\ref{app:asymptotic}) \\
\(q\) & auxiliary budget constant, fixed with \(2<q<Q\); the schedule takes \(s\approx q\log R/\log\log R\) \\
\(\theta_1\) & interval exponent, \(m_1=\floor{(\log R)^{\theta_1}}\), fixed with \(\tfrac12+\tfrac1q<\theta_1<1\) \\
\(c_M\) & the explicit Mertens deficit \((\log4+4)+(2+E_1)\le10\) in the interval floor \\
\(\theta\) & the closing paragraph of Section~\ref{subsec:scales} only: variable exponent \(m_0=(\log R)^{\theta}\) in the optimality discussion \\
\end{tabular}
\end{center}

\subsection*{Section-local symbols}

\begin{center}
\begin{tabular}{@{}lp{0.76\textwidth}@{}}
\(k_R\) & Section~\ref{subsec:cc-log} only: the Cilleruelo--C\'ordoba baseline parameter \(\ceil{\log R/2}\) \\
\(J_C\) & Section~\ref{subsec:cc-log} only: the number of subarcs \(\ceil{Ce/\sqrt2}\) in the baseline argument \\
\(\alpha_j,\gamma_j,\Phi_j\) & Section~\ref{subsec:cc-log} only: the split part \(\prod_j p_j^{\alpha_j}\) of \(R^2\), the exponents \(\gamma_j\) of the angular parameterization, and the angle \(\Phi_j\) of the Gaussian prime above \(p_j\) \\
\(n,\rho,f_b\) & Appendix~\ref{app:missing} collision proof: \(n=2s+1\), the coordinatewise reduction map modulo \(p\), and the fiber counts over residue classes \(b\) \\
\(u\) & Appendix~\ref{app:missing} divisibility proof: the row-scaling prefactor \(\prod_j z_j^s\) that is cancelled
\end{tabular}
\end{center}

\begin{thebibliography}{99}

\bibitem{BombieriGranvillePintz1992}
E. Bombieri, A. Granville and J. Pintz,
\newblock Squares in arithmetic progressions,
\newblock \emph{Duke Mathematical Journal} \textbf{66} (1992), no. 3, 369--385.

\bibitem{BombieriPila1989}
E. Bombieri and J. Pila,
\newblock The number of integral points on arcs and ovals,
\newblock \emph{Duke Mathematical Journal} \textbf{59} (1989), no. 2, 337--357.

\bibitem{BombieriZannier2002}
E. Bombieri and U. Zannier,
\newblock A note on squares in arithmetic progressions, II,
\newblock \emph{Atti della Accademia Nazionale dei Lincei. Rendiconti Lincei.
Matematica e Applicazioni (9)} \textbf{13} (2002), no. 2, 69--75.

\bibitem{BourgainRudnick2011}
J. Bourgain and Z. Rudnick,
\newblock On the geometry of the nodal lines of eigenfunctions of the
two-dimensional torus,
\newblock \emph{Annales Henri Poincar\'e} \textbf{12} (2011), no. 6,
1027--1053.

\bibitem{Chan2015}
T. H. Chan,
\newblock Factors of almost squares and lattice points on circles,
\newblock \emph{International Journal of Number Theory} \textbf{11} (2015),
no. 5, 1701--1708.

\bibitem{Cilleruelo1991}
J. Cilleruelo,
\newblock Arcs containing no three lattice points,
\newblock \emph{Acta Arithmetica} \textbf{59} (1991), no. 1, 87--90.

\bibitem{CillerueloCordoba1992}
J. Cilleruelo and A. C\'ordoba,
\newblock Trigonometric polynomials and lattice points,
\newblock \emph{Proceedings of the American Mathematical Society}
\textbf{115} (1992), no. 4, 899--905.

\bibitem{CillerueloGranville2007}
J. Cilleruelo and A. Granville,
\newblock Lattice points on circles, squares in arithmetic progressions and
sumsets of squares,
\newblock in \emph{Additive Combinatorics}, CRM Proceedings \& Lecture Notes
\textbf{43}, American Mathematical Society, Providence, RI, 2007, 241--262.

\bibitem{CillerueloGranville2009}
J. Cilleruelo and A. Granville,
\newblock Close lattice points on circles,
\newblock \emph{Canadian Journal of Mathematics} \textbf{61} (2009), no. 6,
1214--1238.

\bibitem{HeathBrown2002}
D. R. Heath-Brown,
\newblock The density of rational points on curves and surfaces,
\newblock \emph{Annals of Mathematics (2)} \textbf{155} (2002), no. 2,
553--598.

\bibitem{Jarnik1926}
V. Jarn\'{\i}k,
\newblock \"Uber die Gitterpunkte auf konvexen Kurven,
\newblock \emph{Mathematische Zeitschrift} \textbf{24} (1926), no. 1,
500--518.

\bibitem{KonyaginSteger1994}
S. V. Konyagin and T. Steger,
\newblock On polynomial congruences,
\newblock \emph{Mathematical Notes} \textbf{55} (1994), no. 6, 596--600.

\bibitem{KurlbergLester2025}
P. Kurlberg and S. Lester,
\newblock Poisson spacing statistics for lattice points on circles,
\newblock \emph{Inventiones Mathematicae} \textbf{240} (2025), no. 2,
537--585.

\bibitem{Mertens1874}
F. Mertens,
\newblock Ein Beitrag zur analytischen Zahlentheorie,
\newblock \emph{Journal f\"ur die reine und angewandte Mathematik} \textbf{78} (1874),
46--62.

\bibitem{Mathlib2020}
The mathlib Community,
\newblock The Lean mathematical library,
\newblock in \emph{Proceedings of the 9th ACM SIGPLAN International Conference
on Certified Programs and Proofs}, CPP 2020, 367--381.

\bibitem{deMouraUllrich2021}
L. de Moura and S. Ullrich,
\newblock The Lean 4 theorem prover and programming language,
\newblock in \emph{Automated Deduction -- CADE 28}, Lecture Notes in Computer
Science \textbf{12699}, Springer, 2021, 625--635.

\bibitem{PNTplus}
A. Kontorovich, T. Tao, et al.,
\newblock The PrimeNumberTheoremAnd project,
\newblock \url{https://github.com/AlexKontorovich/PrimeNumberTheoremAnd}, accessed 2026.

\bibitem{Ramana2007}
D. S. Ramana,
\newblock Arithmetical applications of an identity for the Vandermonde
determinant,
\newblock \emph{Acta Arithmetica} \textbf{130} (2007), no. 4, 351--359.

\bibitem{Rudin1960}
W. Rudin,
\newblock Trigonometric series with gaps,
\newblock \emph{Journal of Mathematics and Mechanics} \textbf{9} (1960),
203--227.

\bibitem{Salberger2023}
P. Salberger,
\newblock Counting rational points on projective varieties,
\newblock \emph{Proceedings of the London Mathematical Society (3)}
\textbf{126} (2023), no. 4, 1092--1133.

\bibitem{Szemeredi1974}
E. Szemer\'edi,
\newblock The number of squares in an arithmetic progression,
\newblock \emph{Studia Scientiarum Mathematicarum Hungarica} \textbf{9} (1974),
417.

\bibitem{Temur2020}
F. Temur,
\newblock Discrete fractional integrals, lattice points on short arcs, and
diophantine approximation,
\newblock arXiv:2012.10784 (2020).

\bibitem{Walsh2015}
M. N. Walsh,
\newblock Bounded rational points on curves,
\newblock \emph{International Mathematics Research Notices} \textbf{2015},
no. 14, 5644--5658.

\end{thebibliography}

\end{document}
